% Compile twice: pdflatex -interaction=nonstopmode technical_note_engineering.tex
% Editable reconstruction of technical_note_VE.pdf with audited revisions.
% Computer Modern is deliberately left as the default LaTeX font.
\documentclass[10pt,a4paper]{article}
\usepackage{amsmath,amssymb,amsfonts,booktabs}
\usepackage[margin=1in]{geometry}
\usepackage{hyperref}
\hypersetup{pdftitle={Photovoltaic Layout and Shading Simulator Formulas},
 pdfauthor={Alexandre Littardi},hidelinks}
\newcommand{\src}[1]{\par\noindent{\small\ttfamily Source: \detokenize{#1}}\par}
\newcommand{\clip}{\operatorname{clip}}
\newcommand{\remark}[1]{\par\noindent\textbf{Remark.} #1\par}
\newcommand{\limitation}[1]{\par\noindent\textbf{Model limitation.} #1\par}
\title{Models and Formulas of the Photovoltaic\\Layout and Shading Simulator\\
 \large Mathematical Appendix --- Traceability of Source Code Formulas}
\author{Alexandre Littardi}
\date{Sept 25, 2026\\[4pt]\small Revised against application code: Sept 28, 2026}
\begin{document}
\maketitle
\begin{abstract}
This appendix retraces the calculation models implemented in the Python/Tkinter photovoltaic
layout and shading application. It retains the typography, mathematical notation and section
structure of the original engineering note, while updating the formulas for dated solar
simulations, two-pole DC routes, electrical checks and hourly PV/BESS operation.
All values produced by these models are estimates subject to the stated input and validation
conditions.
\end{abstract}
\newpage
\begingroup\small\tableofcontents\endgroup
\newpage
\section{Notations}
\begin{table}[h]
\centering\small
\begin{tabular}{@{}llp{8.0cm}@{}}
\toprule
Symbol & Unit & Meaning\\ \midrule
$\phi,\lambda$ & $^\circ$ & Latitude north-positive; longitude east-positive\\
$n,N_y$ & day & Day number; 365 or 366 days in calculation year\\
$h,u$ & h & Local decimal clock hour; UTC offset\\
$\gamma,\delta$ & rad & Fractional year angle; declination\\
$E_t$ & min & Equation of time\\
$\omega,\alpha,A_s$ & $^\circ$ & Hour angle; solar elevation; azimuth from north clockwise\\
$\beta,\gamma_p,\theta_i$ & $^\circ$ & Plane tilt; plane azimuth; incidence angle\\
$G_{\rm DNI},G_{\rm POA}$ & W/m$^2$ & Direct normal and plane-of-array irradiance\\
$T_c,T_{\rm amb}$ & $^\circ$C & Cell and ambient temperatures\\
$P_{\rm STC},\mu_P$ & W,\ \%/$^\circ$C & STC module power; power temperature coefficient\\
$f_s,o$ & -- & Shaded fraction; obstacle opacity, both in $[0,1]$\\
$L_A,L_B,S,\rho$ & m,m,mm$^2$,$\Omega$\,mm$^2$/m & Positive and negative cable lengths; section; resistivity\\
$C_{\rm nom},Q$ & kWh & Nominal battery capacity; stored energy\\
$E_{\rm PV},E_{\rm load}$ & kWh & Hourly AC production and refrigerator demand\\
\bottomrule
\end{tabular}
\caption{Main symbols; additional symbols are defined at their first occurrence.}
\end{table}
Trigonometric relations use radians internally; displayed angles are in degrees unless specified.
Electrical energy is in Wh in the shading simulation and kWh in the hourly AC balance.

\section{Solar Position}
\src{mixins/shadow_energy.py: _compute_solar_position}
The position model uses a fractional-year approximation for declination and equation of time.
The preview with day/month alone uses the representative year 2024. Dated shadow and annual
simulations use the actual year; the annual Europe/Rome profile uses the UTC offset of each
individual local hour, including daylight-saving transitions.
\subsection{Fractional Year Angle}
\begin{equation}
\gamma=\frac{2\pi}{N_y}\left(n-1+\frac{h-12}{24}\right),\qquad N_y\in\{365,366\}.
\end{equation}
\subsection{Equation of Time}
\begin{equation}
\begin{split}
E_t=229.18(&0.000075+0.001868\cos\gamma-0.032077\sin\gamma\\
           &-0.014615\cos(2\gamma)-0.040849\sin(2\gamma))\quad[\mathrm{min}].
\end{split}
\end{equation}
\subsection{Solar Declination}
\begin{equation}
\begin{split}
\delta={}&0.006918-0.399912\cos\gamma+0.070257\sin\gamma\\
&-0.006758\cos(2\gamma)+0.000907\sin(2\gamma)\\
&-0.002697\cos(3\gamma)+0.001480\sin(3\gamma).
\end{split}
\end{equation}
\subsection{True Solar Time and Hour Angle}
\begin{align}
t_{\rm offset}&=E_t+4\lambda-60u &&[\mathrm{min}],\\
t_{\rm solar}&=60h+t_{\rm offset} &&[\mathrm{min}],\\
\omega&=\frac{t_{\rm solar}}{4}-180 &&[^\circ].
\end{align}
\subsection{Geometric Solar Elevation}
\begin{equation}
\alpha_{\rm geo}=\arcsin\!\left(\clip_{[-1,1]}
(\sin\phi\sin\delta+\cos\phi\cos\delta\cos\omega)\right).
\end{equation}
\subsection{Atmospheric Refraction Correction}
\begin{equation}
\Delta\alpha_{\rm ref}=
\begin{cases}
\displaystyle \frac{1.02}
{\tan\!\left(\alpha_{\rm geo}+
    \frac{10.3}{\alpha_{\rm geo}+5.11}\right)}
&\alpha_{\rm geo}>-0.85^\circ,\\
0&\text{otherwise}
\end{cases}
\quad[\mathrm{arcmin}],
\end{equation}
\begin{equation}
\alpha=\alpha_{\rm geo}+\frac{\Delta\alpha_{\rm ref}}{60}.
\end{equation}
The constants and angles inside the refraction tangent are expressed in degrees;
the implementation converts them to radians before calling trigonometric functions.
\subsection{Solar Azimuth}
\begin{equation}
c_A=\clip_{[-1,1]}\left(
\frac{\sin\delta-\sin\alpha_{\rm geo}\sin\phi}
{\cos\alpha_{\rm geo}\cos\phi}\right),
\end{equation}
\begin{equation}
A_s=\begin{cases}360^\circ-\arccos(c_A)&\omega>0,\\
\arccos(c_A)&\omega\leq0.\end{cases}
\end{equation}
\limitation{At a singular zenith/polar denominator, the implementation applies a
conventional azimuth instead of evaluating the quotient. These position equations are
approximations, not an ephemeris-grade solar-position algorithm.}

\section{Solar Irradiance on Module Plane}
\src{mixins/shadow_energy.py: get_clear_sky_poa_irradiance}
The shadow-study scenario uses a simplified clear-sky sky model:
$G_{\rm sc}=1367$ W/m$^2$, $\tau=0.7$ and ground albedo $0.20$.
Imported historical weather, when supplied for the annual energy simulation, replaces
the clear-sky plane irradiance and supplies hourly ambient temperature for one orientation.
\subsection{Direct Normal Irradiance}
\begin{equation}
m_a=\frac{1}{\sin\alpha},\qquad
G_{\rm DNI}=
\begin{cases}
G_{\rm sc}\,\tau^{m_a^{0.675}}
&\alpha>0.5^\circ\text{ and }\sin\alpha\geq0.01,\\
0&\text{otherwise}.
\end{cases}
\end{equation}
\subsection{Angle of Incidence}
\begin{equation}
c_i=\clip_{[0,1]}\!
\left(\sin\alpha\cos\beta+\cos\alpha\sin\beta\cos(A_s-\gamma_p)\right).
\end{equation}
\subsection{Direct, Diffuse and Reflected Components}
\begin{align}
G_{\rm dir}&=G_{\rm DNI}c_i,\\
G_{\rm h,diff}&=0.12G_{\rm DNI}\max(0,\sin\alpha),\\
G_{\rm diff}&=G_{\rm h,diff}\frac{1+\cos\beta}{2},\\
G_{\rm h}&=G_{\rm DNI}\max(0,\sin\alpha)+G_{\rm h,diff},\\
G_{\rm alb}&=0.20G_{\rm h}\frac{1-\cos\beta}{2}.
\end{align}
\subsection{Total Plane Irradiance}
\begin{equation}
G_{\rm POA}=\max(0,G_{\rm dir}+G_{\rm diff}+G_{\rm alb}).
\end{equation}
\limitation{The diffuse share is a fixed empirical fraction, and albedo is
constant. The clear-sky scenario is not a real-weather production guarantee.
Imported weather is historical input and must represent the selected plane orientation.}

\section{Module Electrical Model}
\src{mixins/shadow_energy.py: estimate_panel_power_w}
\subsection{Effective Irradiance and Cell Temperature}
\begin{equation}
G_{\rm eff}=G_{\rm POA}(1-f_s),\qquad
T_c=T_{\rm amb}+(T_{\rm NOCT}-20)\frac{G_{\rm eff}}{800}.
\end{equation}
$T_{\rm amb}=20^\circ$C is assumed only for the clear-sky shadow study;
imported hourly weather supplies $T_{\rm amb}$ where enabled.
\subsection{Temperature Factor}
\begin{equation}
k_T=\max\left(0,1+\frac{\mu_P}{100}(T_c-25)\right).
\end{equation}
\subsection{Instantaneous Power}
\begin{equation}
P(t)=
\begin{cases}
P_{\rm STC}\dfrac{G_{\rm eff}}{1000}\,k_T&G_{\rm eff}>0,\\
0&\text{otherwise}.
\end{cases}
\end{equation}
\subsection{Energy over One Step}
For midpoint $t_m=(t_0+t_1)/2$ and $\Delta t$ in hours:
\begin{align}
E_{\rm ideal}&=P_{f_s=0}(t_m)\Delta t,\\
E_{\rm real}&=P_{f_s}(t_m)\Delta t,\\
E_{\rm loss}&=\max(0,E_{\rm ideal}-E_{\rm real}).
\end{align}
\limitation{This is a linear irradiance and temperature power estimate. It does
not solve a diode I--V equation, diode bypass behaviour or series mismatch.}

\section{Time Integration of Shading Simulation}
\src{mixins/shadow_simulation.py: run, data_at}
\begin{equation}
E_{\rm total}=\sum_{k=1}^{M}P(t_{m,k})\,\Delta t_k,\qquad
t_{m,k}=\frac{t_{0,k}+t_{1,k}}{2}.
\end{equation}
The time step controls the accuracy near sunrise, sunset and abrupt shadow changes.
\subsection{Aggregation per String}
\src{mixins/shadow_energy.py: summarize_string_group}
For the $M_s$ modules of a string:
\begin{align}
E_{\rm ideal,string}&=\sum_{j=1}^{M_s}E_{\rm ideal,j},
&E_{\rm loss,string}&=\sum_{j=1}^{M_s}E_{\rm loss,j},\\
r_{\rm loss,string}&=100\frac{E_{\rm loss,string}}{E_{\rm ideal,string}},\\
\Delta_{\rm spread}&=\max_j(r_{\rm loss,j})-\min_j(r_{\rm loss,j}).
\end{align}
The last quantity is a shading heterogeneity indicator, not a calculated
series-string electrical mismatch loss.

\section{Cast Shadow Geometry}
\src{mixins/shadow_geometry.py: compute_shadow_geometry}
The simplified obstacle is a vertical mast with a square plan footprint,
projected separately onto each roof-zone elevation.
\subsection{Polygon Area}
\begin{equation}
A=\frac12\left|\sum_{k=1}^{p}(x_ky_{k+1}-x_{k+1}y_k)\right|.
\end{equation}
\subsection{Projection Distance}
\begin{equation}
d_{\rm shadow}=\frac{H_{\rm mast}-H_{\rm roof}}{\tan\alpha},
\quad H_{\rm mast}>H_{\rm roof},\quad \alpha>0.
\end{equation}
The footprint is translated opposite to the sun azimuth, accounting for
the user-defined image north rotation. The convex hull of original and
translated footprint vertices produces the candidate shadow polygon.
\subsection{Convex Hull and Polygon Intersection}
The convex hull is computed with Andrew's monotone-chain procedure, with
complexity $O(p\log p)$. Intersection with the module rectangle uses successive
Sutherland--Hodgman clipping against each rectangle edge.
\subsection{Shaded Module Fraction}
\begin{equation}
f_s=\clip_{[0,1]}\left(o\,\frac{A_{\rm intersection}}{A_{\rm module}}\right),
\qquad s_{\%}=100f_s.
\end{equation}
\limitation{A zero or negative height difference, or a sun below the horizon,
produces no modeled projected shadow. The obstacle geometry and optical
transparency remain simplified.}

\section{Panel Grid Layout}
\src{mixins/zone_tools.py: recalculate_zone_grids, zone_rotate_point}
\subsection{Rows and Columns}
\begin{equation}
n_{\rm col}=\left\lfloor\frac Ww\right\rfloor,\qquad
n_{\rm row}=\left\lfloor\frac Hh\right\rfloor.
\end{equation}
These dimensions describe the preliminary rectangular grid before individual
edits, roof setbacks and actual mounting constraints.
\subsection{Alignment}
With $r_x=W-n_{\rm col}w$ and $r_y=H-n_{\rm row}h$:
\begin{equation}
o_x\in\{0,r_x/2,r_x\},\qquad o_y\in\{0,r_y/2,r_y\}.
\end{equation}
\subsection{Rotation about Each Module Centre}
\begin{equation}
\begin{pmatrix}x'\\y'\end{pmatrix}=
\begin{pmatrix}c_x\\c_y\end{pmatrix}+
\begin{pmatrix}\cos\theta&-\sin\theta\\\sin\theta&\cos\theta\end{pmatrix}
\begin{pmatrix}x-c_x\\y-c_y\end{pmatrix}.
\end{equation}

\section{DC Wiring}
\src{mixins/two_pole_cables.py: get_two_pole_route}
\subsection{Two Physical Routes}
For endpoints $P_1=(x_1,y_1)$ and $P_2=(x_2,y_2)$, a direct
orthogonal fallback has length
\begin{equation}
L_{\rm orth}=|x_2-x_1|+|y_2-y_1|.
\end{equation}
The drawn pathway network instead uses shortest paths through a graph
(Dijkstra). The string's first and last terminals each receive a separate
route; modeled vertical roof-to-inverter drops and end reserves are included:
\begin{align}
L_A&=L_{{\rm roof},A}+|H_A-H_{\rm inv}|+R_A,\\
L_B&=L_{{\rm roof},B}+|H_B-H_{\rm inv}|+R_B,\qquad
L_{\rm loop}=L_A+L_B.
\end{align}
\limitation{Routes and risers are preliminary estimates until actual cable
trays, inverter elevation, termination points and inter-module wiring have
been surveyed. Recalculate routes after editing panel positions or pathways.}
\subsection{Voltage Drop and Selected Cross-Section}
\src{mixins/notes_tools.py: calculate_dc_cable_size}
\begin{align}
\Delta U_{\rm DC}&=\frac{I\rho(L_A+L_B)}{S},&
S_{\min}&=\frac{I\rho(L_A+L_B)}{U\varepsilon_{\rm target}/100},\\
S_{\rm norm}&=\min\{s\in\mathcal S_{\rm norm}:s\geq S_{\min}\},&
\varepsilon_{\rm actual}&=100\frac{I\rho(L_A+L_B)}{S_{\rm norm}U}.
\end{align}
Here $\mathcal S_{\rm norm}=\{1.5,2.5,4,6,10,16,25,35,50,70,95,120,150,185,240\}$
mm$^2$. In the general-purpose calculator, a specified one-way length $L$
represents $L_A=L_B=L$, giving $L_{\rm loop}=2L$. The project default
resistivity is $0.017\ \Omega\,\mathrm{mm}^2/\mathrm m$ and is user-editable.
\limitation{Voltage drop alone cannot determine conductor acceptability.
Installation temperature, grouping, ampacity, insulation, DC short-circuit
capability, earthing and protective-device coordination require separate
engineering checks.}

\section{Electrical Summary under STC Conditions}
\src{project_validation.py; detailed_electrical.py}
\subsection{One String of \texorpdfstring{$N$}{N} Modules in Series}
\begin{align}
P_{{\rm dc},s}&=\frac{NP_{\rm STC}}{1000}\ [\mathrm{kWp}],&
V_{{\rm oc},s}&=NV_{\rm oc},&
V_{{\rm mpp},s}&=NV_{\rm mpp},\\
I_{{\rm sc},s}&=I_{\rm sc},&
I_{{\rm mpp},s}&=I_{\rm mpp}.&&
\end{align}
\subsection{Parallel Strings on One MPPT}
\begin{equation}
P_{{\rm dc},m}=\sum_{s\in m}P_{{\rm dc},s},\qquad
I_{{\rm sc},m}=\sum_{s\in m}I_{{\rm sc},s},\qquad
I_{{\rm mpp},m}=\sum_{s\in m}I_{{\rm mpp},s}.
\end{equation}
Strings of unequal series lengths on the same MPPT are flagged for review.
An interval of string voltages does not demonstrate their electrical
compatibility; module technology, operating temperature, currents and
manufacturer input limits must be verified.
\subsection{Cold Open-Circuit Voltage and Inverter Totals}
\begin{equation}
V_{{\rm oc},s}(T_{\min})=NV_{\rm oc,STC}
\left[1+\frac{\mu_{\rm Voc}}{100}(T_{\min}-25)\right],
\end{equation}
\begin{equation}
P_{{\rm dc},{\rm inv}}=\sum_{s\ {\rm assigned}}P_{{\rm dc},s},\qquad
I_{{\rm ac},{\rm nominal}}=
\frac{1000S_{{\rm ac},{\rm kVA}}}{\sqrt3\,U_{\rm LL}}.
\end{equation}
The cold-voltage calculation needs the module coefficient and design minimum
cell temperature. The AC expression assumes three-phase operation and a
specified line-to-line voltage; it is not a breaker setting.
The project arrangement uses three hybrid inverters of nominal 125 kW with
10 MPPT each. Product-specific operating windows and BESS interfaces must
be confirmed from the exact equipment documentation.

\section{Hourly PV, Refrigerator Load and BESS}
\src{self_consumption.py; battery_dispatch.py}
The following AC-side accounting is performed independently for each local
hour ($\Delta t=1$ h). For the PV-only option, set charge and discharge to zero.
\subsection{Direct Use and Storage Dispatch}
\begin{align}
E_{\rm direct}&=\min(E_{\rm PV},E_{\rm load}),\\
E_{\rm surplus}&=E_{\rm PV}-E_{\rm direct},&
E_{\rm deficit}&=E_{\rm load}-E_{\rm direct}.
\end{align}
Let $Q\in[Q_{\min},Q_{\max}]$ be stored energy, and let
$\eta_c=\eta_d=\sqrt{\eta_{\rm roundtrip}}$. At each hour:
\begin{align}
E_{\rm charge}&=\min\!\left(E_{\rm surplus},P_c\Delta t,
                       \frac{Q_{\max}-Q}{\eta_c}\right),\\
Q'&=Q+\eta_c E_{\rm charge},\\
E_{\rm discharge}&=\min\!\left(E_{\rm deficit},P_d\Delta t,
                       (Q'-Q_{\min})\eta_d\right),\\
Q_{\rm next}&=Q'-\frac{E_{\rm discharge}}{\eta_d}.
\end{align}
The default battery is provisionally identified as
\emph{Deye MC-L522-2H3}, with nominal $C_{\rm nom}=522.496$ kWh per
cabinet, $Q_{\min}=0.10 C_{\rm nom}$, $Q_{\max}=0.90 C_{\rm nom}$,
initial $Q=Q_{\min}$ and $\eta_{\rm roundtrip}=0.90$.
The configured limits $P_c=250$ kW and $P_d=150$ kW are \emph{shared
plant limits}: two cabinets double energy capacity but do not double power.
Charging from the grid and resale of discharged battery energy are not modeled.
\subsection{Grid Import, Export and Curtailment}
\begin{align}
E_{\rm grid}&=E_{\rm deficit}-E_{\rm discharge},\\
E_{\rm export}&=\min(E_{\rm surplus}-E_{\rm charge},
                    P_{\rm export}\Delta t),\\
E_{\rm curtailed}&=E_{\rm surplus}-E_{\rm charge}-E_{\rm export},\\
E_{\rm useful\ self}&=E_{\rm direct}+E_{\rm discharge},\\
E_{\rm PV\ retained}&=E_{\rm direct}+E_{\rm charge}.
\end{align}
The last two totals differ because of battery losses and changes in stored
energy; they must not be added. Checks on annual aggregates are
$E_{\rm load}=E_{\rm useful\ self}+E_{\rm grid}$ and
$E_{\rm PV}=E_{\rm PV\ retained}+E_{\rm export}+E_{\rm curtailed}$.
The existing grid contract is recorded as 507 kW, and 300 kW export is
the \emph{requested, not approved} study limit.

\section{Annual Economics and Scenario Comparison}
\src{battery_dispatch.py; mixins/self_consumption_ui.py}
Annual energy quantities are sums of the hourly quantities above. Let
$p_b$ be the purchase price in EUR/kWh and $p_s$ the credit for exported
energy in EUR/kWh. Then for each scenario $j$:
\begin{align}
C_{0}&=p_b E_{{\rm load},\rm annual},\\
C_j&=p_b E_{{\rm grid},j}-p_s E_{{\rm export},j},\\
B_j&=C_0-C_j,\\
T_j&=\begin{cases}K_j/B_j&B_j>0,\\
\text{undefined}&B_j\leq0.\end{cases}
\end{align}
$C_j$ is the \emph{net annual electricity cost} after export credit;
$B_j$ is the annual benefit relative to buying the entire refrigerator
demand. For PV only, one battery and two batteries, illustrative upfront
capital $K_j$ is respectively EUR 250,000, 300,000 and 350,000.
Evaluate $p_b=0.25$ and $0.32$ EUR/kWh and a specified export tariff in
the proposed $0.07$--$0.08$ EUR/kWh range. A baseline of EUR 125,000/year
at EUR 0.25/kWh corresponds to 500,000 kWh/year \emph{only if} the
metered annual demand confirms that figure. The simple payback $T_j$
excludes operating expenses, tax, financing, degradation, battery
replacement and future tariff changes.

\section{Summary of Model Limitations and Required Engineering Checks}
\begin{itemize}
\item The default sky is simplified clear sky. Hourly historical weather can
be imported, but prediction and uncertainty still depend on input quality.
\item Power scaling and string shading do not simulate diode bypass or I--V
curves. The spread indicator is diagnostic only.
\item Inverter count, MPPT allocation, cold $V_{\rm oc}$, hot $V_{\rm mpp}$,
parallel current and any bifacial contribution need actual datasheets.
\item Two-pole DC routes include modeled vertical drops and reserves;
field routing, AC cables, inter-module leads and protection studies remain
to be completed by the responsible designer.
\item BESS power and efficiency are configurable assumptions; the listed
model and interface must be confirmed before procurement. Export at
300 kW remains subject to grid-operator approval.
\end{itemize}
\end{document}
